\documentclass[11pt]{article}
\usepackage[T1]{fontenc}
\usepackage[margin=1in]{geometry}
\usepackage{amsmath,amssymb,amsthm,array}
\usepackage[hidelinks]{hyperref}
\newcommand{\Q}{\mathbb Q}
\newcommand{\C}{\mathbb C}
\newcommand{\R}{\mathbb R}
\newcommand{\kk}{{\overline{\mathbb Q}}}
\DeclareMathOperator{\End}{End}
\DeclareMathOperator{\Hom}{Hom}
\DeclareMathOperator{\Lie}{Lie}
\DeclareMathOperator{\Gal}{Gal}
\DeclareMathOperator{\Tr}{Tr}
\DeclareMathOperator{\Trd}{Trd}
\DeclareMathOperator{\Hg}{Hdg}
\DeclareMathOperator{\SL}{SL}
\DeclareMathOperator{\Sp}{Sp}
\DeclareMathOperator{\SO}{SO}
\newtheorem{theorem}{Theorem}
\newtheorem{lemma}[theorem]{Lemma}
\newtheorem{proposition}[theorem]{Proposition}
\newtheorem{corollary}[theorem]{Corollary}
\theoremstyle{remark}
\newtheorem{remark}[theorem]{Remark}
\title{Actual derived Hodge images in dimension at most six:\\
a universal reduction and an exact certificate}
\author{Campaign 9000032: Engineer submission}
\date{November 5, 2026}
\begin{document}
\maketitle

\begin{abstract}
We give a verification proof of the actual-group independence assertion
\texttt{p3ab:independence} and its classification prerequisites in
\emph{Weil classes and Hodge classes on abelian powers}, October 6, 2026.
The statement concerns every complex abelian variety of dimension at most
six with a unital imaginary-quadratic action, not just a generic period.
The proof uses the full endomorphism commutant, coefficient descent,
Hodge-circle minimality, and the prime-sized Galois orbits that occur at
partially pure embeddings. It treats the exceptional identifications
$B_2=C_2$, $D_2=A_1+A_1$, and $D_3=A_3$ explicitly.
An executable rational-arithmetic certificate enumerates all necessary
numerical cases and all possible subdirect graph partitions after these
reductions. This is a universal proof with a finite algebraic certificate,
not an enumeration of sampled abelian varieties.
It verifies the derived-image condition (c) of the manuscript's tensor
criterion; it does not verify its determinant-supply condition (d), the
Hodge conjecture, or the remaining manuscripts.
\end{abstract}

\noindent\textbf{Certification status.}
This is an Engineer proof submission with a successfully recomputed exact
certificate. Independent Reviewer acceptance, independent attribution of
the cited inputs, and publication-level significance/novelty adjudication
remain pending. No claim of verified-new priority or campaign completion
is made.

\section{Precise statement and imported inputs}

Fix $\kk=\overline{\Q}\subset\C$. Scalar conjugation means changing
algebraic coefficients in the cohomology of a fixed complex variety, not
replacing that variety by a conjugate. For $H=H^1(A,\Q)$, let
$J=\Hg(H)$ be the smallest rational algebraic subgroup containing the
Hodge circle. The circle has weights $+1,-1$ on $H^{1,0},H^{0,1}$.
Put $R=J^{\rm der}$. A summand or map over $\kk$ is
\emph{admissible} if its projector or map preserves the Hodge grading in
every scalar conjugate. A pair tensor is nondegenerate and of type
$(1,1)$ in every scalar conjugate; its symmetry is ordinary tensor
symmetry, before any cohomological Koszul sign.

Write, up to isogeny,
\[
 A\sim\prod_B B^{r_B}
\]
with the $B$ simple and pairwise nonisogenous. A \emph{column} means the
image of a primitive matrix idempotent in the scalar extension of the
\emph{actual full} algebra
\[
 D_B=\End_{\rm Hdg}(H^1(B,\Q))=\End^0(B)^{\rm op}.
\]
It is not a summand obtained merely from the initially supplied quadratic
field. Matrix copies selected by this algebra are identified. In Type IV
we retain one representative from each pair of opposite center embeddings,
and act dually on its opposite.

\begin{theorem}\label{thm:main}
Let $A/\C$ have dimension at most six and admit a unital embedding of an
imaginary quadratic field $E$ into $\End^0(A)$. No simplicity, genericity,
centrality of the action, or hyperbolicity assumption is imposed.
Discard the CM isogeny factors. The remaining actual columns have the
following necessary possibilities, with $g=\dim B$.
\begin{enumerate}
\item In Albert Types I and II their dimensions are $2,4,6$, and the
derived projection is the full symplectic group.
\item In Type III, $g=4$ or $6$, the center is $\Q$, and $r_B=1$.
There is one column, with image $\SO_g$. The two eigenspaces of $E$
are its two matrix copies, each with Hodge dimensions $(g/2,g/2)$.
\item In Type IV the actual algebra is a CM field $K$.
If $E$ does not embed in $K$, then $[K:\Q]=2$, $g=3$, $r_B=2$,
and the column image is $\SL_3$.
If $E$ embeds in $K$, put $[K:\Q]=2f$ and $a=g/f$. Then
\[
 (f,a,g)=(1,g,g),\quad 3\leq g\leq6,
 \qquad\text{or}\qquad
 (f,a,g)\in\{(2,2,4),(2,3,6),(3,2,6)\}.
 \]
Every column has image $\SL_a$, with dual action at the opposite
embedding. If $a=2$, the embedding pairs include both pure and mixed
ones. A pure column has only one circle weight.
\end{enumerate}
After the stated matrix-copy and opposite-embedding identifications,
these images act independently. More precisely, the image of $R_\kk$
on all of $H_\kk$ is
\[
 \prod_i\Sp(V_i)\ \times\ \prod_j\SO(V_j)\
 \times\ \prod_l\SL(V_l),
 \tag{1}\label{eq:product}
\]
with standard action on its own copies, dual action on opposite linear
copies, and trivial action on the CM summand and on other blocks.
The forms and copy identifications are admissible.
\end{theorem}

The list is a \emph{covering list}, not an existence theorem for every
listed endomorphism algebra or signature. This distinction permits the
certificate to overapproximate realizable cases without losing a case.
Rank-two Type IV columns remain \emph{linear} blocks: derived self-duality
alone does not supply a Hodge self-pair.

The proof uses four established structural inputs.
They are recorded as assumptions in the ledger
\path{research/MATH_STATE.json}.
\begin{center}\small
\begin{tabular}{p{.18\textwidth}p{.73\textwidth}}
\texttt{ALBERT} &
The classification of positive endomorphism division algebras:
Chai--Oort \cite[Sections 10.12--10.14]{CO}, with the Type III wording
also checked against \cite[Section (1.1)]{MZ}.\\[3pt]
\texttt{HODGE-GROUP} &
The actual Hodge group is connected reductive, its commutant is the
endomorphism algebra, and it is toric exactly in the CM case:
\cite[Sections (1.2)--(1.3)]{MZ}. The source uses homology, explaining
the opposite algebra in our cohomological convention.\\[3pt]
\texttt{DYNKIN} &
The root-system classification and low-rank identifications:
\cite[Example 23.5, Theorem 23.7, Remark 23.18]{Et}.\\[3pt]
\texttt{MINUSCULE-}\newline\texttt{DIMENSION} &
Highest-weight dimensions, minuscule weights and their Weyl orbits:
\cite[Proposition 26.8, Definition 30.1, Lemma 30.2,
Corollary 30.5, Sections 27.2 and 30.4]{Et}.
\end{tabular}
\end{center}
The exact cited versions and passages have been obtained and read;
the PDFs, extracted text and short locator-specific excerpts are in
\texttt{literature/}. Independent citation attribution is not replaced
by the Engineer's own reading. No CM Hodge-conjecture theorem, geometric
Weil-plane theorem, or assertion of the audited manuscript is imported
as a premise of Theorem~\ref{thm:main}.

\paragraph{Known scope and comparison.}
Moonen--Zarhin \cite[Theorems (0.1)--(0.2)]{MZ} already treat all abelian
varieties through dimension five. Their Section (3.1) describes graph
links; Lemmas (3.3)--(3.4) impose explicit non-compact-type and
representation-uniqueness hypotheses. We do not silently apply those
lemmas to partially pure columns. The additional work here is an
all-case bridge for this dimension-six quadratic-action problem and
a reproducible algebraic certificate. A bounded literature lookup
located this predecessor but did not establish publication novelty of
the dimension-six formulation. In particular, a search result that
merely points back to the manuscript under audit cannot verify it.
The preexisting campaign ledger contained no verification of this step;
the present contribution is a universal proof reduction, rather than
a larger finite sample.

\section{Polarization, actual columns, and rational descent}

\subsection{The polarization is an actual one}

Write $E=\Q(u)$, $u^2=-d$, $d>0$. If $e_0$ is any rational Riemann
form, set
\[
 e(x,y)=e_0(x,y)+d^{-1}e_0(ux,uy).
 \tag{2}\label{eq:average}
\]
Because $u$ is a complex-linear invertible endomorphism, both summands
are positive Riemann forms. Directly,
$e(ux,y)=-e(x,uy)$; thus $e$ is compatible with $E$.
A positive integer multiple clears denominators. In the manuscript's
homological convention,
\[
 h_E(x,y)=e(x,uy)+u e(x,y)
 \tag{3}
\]
is conjugate-linear first, linear second, and Hermitian. Indeed
$h_E(ux,y)=-u h_E(x,y)$, $h_E(x,uy)=u h_E(x,y)$, and
$h_E(y,x)=\overline{h_E(x,y)}$ follow from $u^2=-d$ and alternation.
The trace identity is
\[
 e(x,y)=\Tr_{E/\Q}\bigl((2u)^{-1}h_E(x,y)\bigr).
\]
The independence proof needs a polarization and its actual Rosati
involution, but no assertion that this Hermitian form is hyperbolic.
We consequently impose no unverified signature or rational-discriminant
restriction.

\begin{lemma}[Coefficient descent]\label{lem:descent}
Let $T$ be a rational vector space, $P\subset T_\C$ a complex-linear
subspace, and $t\in T_\kk$. If $\gamma t\in P$ for every
$\gamma\in\Gal(\kk/\Q)$, then
$t\in(T\cap P)\otimes_\Q\kk$.
\end{lemma}
\begin{proof}
Choose a finite Galois extension $L/\Q$ containing the coefficients
of $t$, and a basis $b_i$ with trace-dual basis $b_i^*$.
The rational coefficient of $b_i$ in $t$ is
\[
 t_i=\Tr_{L/\Q}(b_i^*t)
     =\sum_{\sigma\in\Gal(L/\Q)}\sigma(b_i^*)\,\sigma(t).
 \]
Every summand lies in $P$, so $t_i\in T\cap P$. This proves the claim.
Extending each $\sigma$ to $\kk$ justifies the use of the hypothesis.
\end{proof}

In particular, an algebraic-coefficient operator whose every conjugate
is a Hodge endomorphism lies in the scalar extension of the \emph{full
rational} Hodge commutant. It is not sufficient to commute with the
circle at one embedding. Scalar conjugation here is never assumed to
commute with complex conjugation on all of $\kk$.

\begin{lemma}[Actual Albert columns]\label{lem:columns}
For a simple non-CM $B$, write $D=D_B$, $Z=Z(D)$ and $[D:Z]=s^2$.
Then
\[
 D_\kk=\prod_{\tau:Z\hookrightarrow\kk}M_s(\kk),\qquad
 H^1(B,\kk)=\bigoplus_\tau W_\tau\otimes U_\tau,\quad \dim U_\tau=s.
 \tag{4}\label{eq:columns}
\]
The $W_\tau$ are irreducible for both $J_B$ and $R_B$, distinct as
full-$J_B$ modules, and matrix units give admissible copy
identifications. Their necessary data are
\[
\begin{array}{c|c|c|c}
\text{type}&\text{divisibility}&\dim W_\tau&\text{pair}\\ \hline
\mathrm I&f\mid g&2g/f&\text{alternating self-pair}\\
\mathrm{II}&2f\mid g&g/f&\text{alternating self-pair}\\
\mathrm{III}&2f\mid g&g/f&\text{symmetric self-pair}\\
\mathrm{IV}&fb^2\mid g&g/(fb)&\text{opposite-embedding pair}
\end{array}
\tag{5}\label{eq:albert}
\]
Here $[Z:\Q]=f$ in the first three types, while $Z=K$ is CM of
degree $2f$ and $[D:K]=b^2$ in Type IV.
\end{lemma}
\begin{proof}
The Hodge-group input identifies the commutant with the actual division
algebra. Complete reducibility and the centralizer of a direct sum of
irreducibles give (4); centralizers commute with scalar extension
because their defining equations are linear. The connected center of
$J_B$ acts by scalars on each $W_\tau$, and $Z(J_B)^0R_B\to J_B$ is
surjective. Thus a nontrivial $R_B$-submodule would be a
$J_B$-submodule. All projectors and matrix units lie in $D_\kk$, so
their conjugates remain Hodge maps.

Since $H^1(B,\Q)$ is a vector space over $D$, dimensions give
$2f\mid g$ in the quaternionic cases and $fb^2\mid g$ in Type IV.
Dividing each central block by its matrix multiplicity gives the
column dimensions in (5).

For completeness, the pairing signs do not require a generic-group
assumption. Let $\psi$ be the rational polarization pairing. On a
first-kind central block, its adjoint identity with $M_s(\kk)$ gives
\[
 \psi|_{W\otimes U}=\beta_W\otimes b_U.
 \tag{6}
\]
To see this, fix the $W$ arguments. A bilinear form on $U$ with
$c(Tx,y)=c(x,T^\dagger y)$ differs from the form defining
$\dagger$ by an operator commuting with every matrix, hence by a
scalar. Nondegeneracy gives nondegenerate factors in (6).
Their symmetry signs are opposite.

In Type I, $b_U$ is symmetric. In Type II, the positive involution on
$M_2(\R)$ is orthogonal: a symplectic-adjoint involution would give
$e^\dagger=1-e$ for a rank-one idempotent and violate positivity of
$\Tr(ee^\dagger)$. In Type III, positivity on the Hamilton quaternions
forces canonical conjugation. An orthogonal involution has a nonzero
fixed pure quaternion $v$ and would give
$\Trd(vv^\dagger)=\Trd(v^2)<0$. Canonical conjugation becomes
symplectic-adjoint on $M_2(\kk)$. These facts give exactly the three
self-pair signs in (5). In particular the Type I dimension is even,
which gives $f\mid g$.

This argument supplies forms on \emph{actual chosen columns}. For
Types I--II restrict $\psi$ to a copy selected by a nonisotropic line
for $b_U$ and transport by matrix units. For Type III use two
multiplicity lines with nonzero $b_U$-pairing, then identify their
copies by a matrix unit. Restriction twice to a single multiplicity
line would be zero and is not the construction. All these operations
and their scalar conjugates are Hodge maps, and a nonzero determinant
remains nonzero on conjugating. The inverse tensors are consequently
pair tensors of the required signs.

In Type IV the positive involution restricts to CM conjugation on $K$.
The rational polarization pairs opposite central blocks perfectly;
factoring the matrix multiplicities gives perfect opposite-column
pairs, again transported to the chosen copies. These identify opposite
columns as derived duals. Oppositeness is transported by the rational
pairing, not by an unjustified commutation of arbitrary Galois
automorphisms with complex conjugation.
\end{proof}

\section{Two weights and the exhaustive small representations}

\begin{lemma}[Minimality at all scalar embeddings]\label{lem:minuscule}
Every nontrivial simple-ideal constituent of an actual column is
minuscule, including a constituent whose column is pure at the current
embedding.
\end{lemma}
\begin{proof}
Normalize the complexified circle differential to have eigenvalues
$+1,-1$. A Galois orbit of simple ideals of $\Lie(R_B)_\kk$ descends
to a rational adjoint factor. The differential cannot vanish on that
entire orbit: the kernel of the corresponding rational adjoint action
would be a proper rational subgroup containing the connected Hodge
circle, contrary to minimality.

Fix an acting ideal and its irreducible representation in the tensor
decomposition of a column. Move both by one scalar conjugation so
that this ideal has nonzero circle projection $x$. The transported
column is still admissible, so still has at most the two circle
weights. Fixing eigenvectors in all other tensor factors shows that
$x$ has at most two eigenvalues on this factor. It is semisimple,
traceless and nonzero: a nonzero representation of a simple Lie
algebra is faithful. It therefore has at least two eigenvalues.

Choose a Cartan containing $x$. If a highest weight $\lambda$ were
not minuscule, some positive root $\alpha$ would have
$\langle\lambda,\alpha^\vee\rangle=m\geq2$. For a highest vector
$v$ and a root $\mathfrak{sl}_2$ triple, the identities
$efv=mv$ and $ef^2v=2(m-1)fv$ show that $v,fv,f^2v$ are nonzero.
Thus the weights include
$\lambda,\lambda-\alpha,\lambda-2\alpha$. The Weyl orbit of a root
spans the Cartan dual of an irreducible root system, so a Weyl
conjugate $x'$ of $x$ has $\alpha(x')\ne0$. Those three weights give
three distinct eigenvalues on $x'$, hence on $x$, a contradiction.
Transport the conclusion back to the original ideal. The passage to
an active conjugate, rather than an assumption that the original
column is mixed, is essential.
\end{proof}

All relevant columns have dimension at most six: Type I has even
$r_B$ by the quadratic-action argument below, hence $g\leq3$;
the other types already have column dimension at most $g\leq6$.
A simple Lie algebra acting faithfully on an $n$-space has a Cartan
inside the traceless diagonal matrices after diagonalization, so its
rank is at most $n-1\leq5$. The root-system list is therefore
provably finite before any computation.

The certificate constructs every irreducible Cartan type of rank at
most five: $A_1,\ldots,A_5$, $B_2,\ldots,B_5$,
$C_2,\ldots,C_5$, $D_4,D_5,F_4,G_2$. Reflections generate all roots
and coroots. For a fundamental weight $\omega_i$, the minuscule test
is that every positive coroot has $i$th coefficient at most one.
The dimension is recomputed as
\[
 \dim L_{\omega_i}
 =\prod_{\beta^\vee=\sum_j c_j\alpha_j^\vee>0}
       \frac{\sum_j c_j+c_i}{\sum_j c_j}.
 \tag{7}
\]
Its Weyl-orbit cardinality is checked independently. The resulting
dimension-at-most-six table is
\[
\begin{array}{c|l}
2&A_1\text{ standard, alternating}\\
3&A_2\text{ standard and dual}\\
4&A_3\text{ standard and dual};\quad C_2=B_2\text{ standard/spin,
alternating}\\
5&A_4\text{ standard and dual}\\
6&A_5\text{ standard and dual};\quad C_3\text{ standard,
alternating};\quad A_3\text{ on }\Lambda^2\kk^4,\text{ symmetric}.
\end{array}
\tag{8}\label{eq:small}
\]
The unmarked entries are nonselfdual. Self-duality is also checked
by testing whether $-\omega_i$ belongs to its Weyl orbit.
There are thirteen entries if the $B_2/C_2$ alias is counted twice,
and twelve distinct module types up to that identification.

An irreducible module of a semisimple Lie algebra is the tensor
product of irreducibles for its acting simple ideals. The only
nontrivial products of dimensions at most six are $2\cdot2$ and
$2\cdot3$. The former is a tensor product of two alternating
$A_1$ standards and is symmetric; the latter is nonselfdual and
has different simple types $A_1,A_2$. Three nontrivial factors
already have dimension at least eight.

The executable checks do not infer form signs from dimensions.
They solve $X^tB+BX=0$ and $XT-TX=0$ over $\Q$ in explicit
standard, symplectic, tensor, and exterior-square models.
For $\kk^2\otimes\kk^2$ they give a unique nondegenerate symmetric
form, scalar commutant, and image dimension six. For
$\Lambda^2\kk^4$ they give a unique nondegenerate symmetric form,
scalar commutant, and image dimension fifteen; all bracket
identities of the exterior-square representation are checked.
These are precisely $\mathfrak{so}_4$ and $\mathfrak{so}_6$,
since the respective orthogonal Lie algebras have those dimensions.
The Cartan permutations for $B_2=C_2$ and $D_3=A_3$ are also
computed. No $\SO_3$ or $\SO_5$ vector representation is
minuscule here, and $D_4$ triality representations are too large.

\begin{corollary}\label{cor:selfpaired}
An alternating self-paired column of dimension $2,4,6$ has full
symplectic image. A symmetric self-paired column cannot have dimension
two; in dimensions four and six it has full $\SO_4$ and $\SO_6$
image, respectively.
\end{corollary}
\begin{proof}
Self-duality of an irreducible tensor product forces self-duality
of each simple-ideal factor, by restriction to that ideal. Apply
(8), its two tensor possibilities, and the computed form signs.
Alternatively a connected subgroup of $\SO_2$ is toric and cannot
act irreducibly on a two-dimensional space over $\kk$.
Connectedness upgrades the Lie images to the stated matrix groups.
\end{proof}

\begin{lemma}[Quadratic full center]\label{lem:quadratic}
If the actual $D_B=K$ is imaginary quadratic and $B$ is non-CM,
then a column of dimension at most six has full special linear
derived image. Its dimension is at least three.
\end{lemma}
\begin{proof}
The two columns $W,\overline W$ are both mixed. Otherwise the circle
lies in the rational torus of $K$-scalars, making $J_B$ toric.
If $W$ were derived-selfdual, the opposite pairing would give a
one-dimensional space $\Hom_{R_B}(W,\overline W)$ of isomorphisms.
Normality of $R_B$ makes this line $J_B$-stable. A nonzero $t$ on
it satisfies $h(z)t h(z)^{-1}=z^m t$ for an integer $m$.
Both source and target spectra are $\{-1,1\}$, so translation by
$m$ preserves this set and $m=0$. Every scalar conjugate has the
same property. Lemma~\ref{lem:descent}, after extending by the source
projector, places $t$ in $D_{B,\kk}$. This contradicts preservation
of the two distinct central embeddings.

Every simple ideal of $\Lie(R_B)_\kk$ acts on $W$, because it acts
dually on $\overline W$ and the total representation is faithful.
Each Galois orbit of ideals has an active circle projection by
Lemma~\ref{lem:minuscule}. Two active ideals on the same tensor
product would give at least three eigenvalues: the sum of two
nonconstant finite spectra in characteristic zero has at least
three elements. Thus there is only one Galois orbit of acting ideals,
all of the same Lie type. The only composite nonselfdual possibility
in dimension at most six is $2\otimes3$, of distinct types $A_1,A_2$,
so it cannot occur. The nonselfdual simple entries in (8) are exactly
the standard or dual $\SL_n$ modules, $3\leq n\leq6$.
\end{proof}

\section{Coverage of all quadratic-action cases}

The isotypic projectors are central in $\End^0(A)$, so the $E$-action
restricts to a unital action on every $B^{r_B}$. A unital map from
a field into a nonzero algebra is injective.

In Type I, the center $F$ is totally real. The algebra $F\otimes_\Q E$
is a quadratic field over $F$ and acts on $F^{r_B}$. Hence $r_B$ is
even, $g\leq3$, and (5) gives alternating dimensions $2,4,6$.
In Type II, $2f\mid g\leq6$ gives the same dimensions.
Corollary~\ref{cor:selfpaired} gives their projections.

In Type III that corollary eliminates column dimension two.
The remaining $g/f=4$ or $6$, together with $gr_B\leq6$ and
$2f\mid g$, forces $f=r_B=1$ and $g=4$ or $6$.
Now $E\hookrightarrow D_B$ splits to
$\kk\times\kk\hookrightarrow M_2(\kk)$. Its two nonzero complementary
idempotents both have rank one. Thus its eigenspaces are the two
matrix copies of the column, and the self-pair gives equal Hodge
dimensions $g/2,g/2$. This uses the actual quaternionic
multiplicity space, not an assumed central $E$-action.

For Type IV the following elementary exclusion is useful.
\begin{lemma}[Rank-two and pure columns]\label{lem:pure}
If all columns of a simple factor are pure, that factor is CM.
A Type IV rank-two column cannot have every scalar conjugate mixed.
Consequently a non-CM rank-two Type IV factor has both pure and
mixed embedding pairs.
\end{lemma}
\begin{proof}
If every column is pure, the circle acts scalarly on each entire
central block, including its matrix multiplicity. It lies in the
rational torus of center scalars (with the opposite-pair norm
condition), so minimality makes the Hodge group toric.

If a rank-two column $W$ is mixed in every conjugate, any nonzero
element of $\Lambda^2W$ is a nondegenerate alternating pair tensor
of type $(1,1)$ in every conjugate. Compose its inverse with the
polarization pair $W^*\to\overline W$. The Hodge shifts cancel,
giving an admissible isomorphism $W\to\overline W$.
Coefficient descent puts its extension in the full commutant,
contradicting the distinct center embeddings. This argument works
with matrix multiplicities as well.
\end{proof}

If $b>1$, the divisibility $fb^2\mid g\leq6$ leaves only
$(g,f,b,\dim W)=(4,1,2,2)$. There is only one opposite pair,
which is either pure or mixed in every conjugate; both are
excluded for a non-CM Type IV factor by Lemma~\ref{lem:pure}.
Thus $D_B=K$ is a field. Rank one is likewise toric.

The central $K$-action commutes with the given action on
$K^{r_B}$, so $K\otimes_\Q E$ acts. If $E$ does not embed in
$K$, this is a quadratic field over $K$, requiring even $r_B$
and $g\leq3$. Rank one and the quadratic rank-two case having
been excluded, the only non-CM possibility is
$(f,g,r_B,a)=(1,3,2,3)$, covered by Lemma~\ref{lem:quadratic}.

If $E$ embeds in $K$, then $K\otimes_\Q E=K\times K$.
Its two idempotents on $K^{r_B}$ have arbitrary ranks
$r_+,r_-\geq0$, $r_++r_-=r_B$; either may be zero.
This explicitly includes noncentral actions such as
$a\mapsto\operatorname{diag}(a,\bar a)$.
Conjugating those idempotents by a $K$-matrix only reallocates
matrix copies to $E$-eigenspaces and does not change the derived
columns or their admissible identifications.
Now $fa=g\leq6$, $a\geq2$, and the quadratic rank-two exclusion
give exactly the Type IV list in Theorem~\ref{thm:main}.
For $f=1$ use Lemma~\ref{lem:quadratic}. For $f>1$, $a=2$ or
$3$; irreducibility and (8) force a single simple image
$\SL_2$ or $\SL_3$. A higher-dimensional column cannot be
derived-trivial, since the center would act by scalars on an
irreducible space of dimension greater than one. This remains
valid at a pure embedding. The rank-two assertion follows from
Lemma~\ref{lem:pure}.

Finally the restriction $J_A\to J_B$ is surjective, not just a
generic-period assertion. Its rational image contains the circle of
$B$, hence $J_B$ by minimality. Conversely the inverse image of
$J_B$ contains the circle of $A$, hence is all of $J_A$.
Surjectivity of derived groups follows for a surjection of connected
reductive groups. Thus all the projections computed above are
projections of the \emph{actual} $R_A$.

\section{Universal exclusion of graph links}

\begin{lemma}[Subdirect partitions]\label{lem:graphs}
Split every $\mathfrak{so}_4$ target into its two simple ideals.
A semisimple Lie algebra projecting onto all the resulting simple
targets is a product of graph-diagonal simple ideals. Failure to be
the full product therefore identifies at least two isomorphic
target ideals. The equivalence relation of linked targets is
Galois-invariant for the actual rational Hodge group.
\end{lemma}
\begin{proof}
Decompose the source into simple ideals. The image of each in a
simple target is an ideal, hence either zero or the whole target;
in the latter case the map is an isomorphism. Two different source
ideals cannot both surject to that target, because they commute.
Thus each target has a unique source ideal, and grouping targets
with the same source gives exactly the assertion. Rationality
transports source ideals, target ideals and their projection maps
under scalar conjugation.
\end{proof}

\subsection{Ordinary mixed columns}

Suppose two linked target modules are mixed in every scalar
conjugate. Apart from the $\SO_4$ factors and the
$\SO_6/\SL_4$ distinction below, table (8) says that an isomorphism
of their target Lie algebras identifies the modules, up to the
standard/dual choice for $\SL$. Replace one column by its opposite
when necessary. The derived intertwiner equations are over $\kk$;
there is a nonzero intertwiner $t$ over $\kk$, and its space is a
line by irreducibility. As in Lemma~\ref{lem:quadratic}, normality
makes it a full-$J_A$ character line. Its circle character translates
$\{-1,1\}$ into $\{-1,1\}$, so is trivial. Applying exactly the
same argument to each conjugate, not conjugating the circle itself,
and then applying Lemma~\ref{lem:descent}, gives an element of
\[
 \End_{\rm Hdg}(H^1(A,\Q))\otimes_\Q\kk.
 \tag{9}
\]
This algebra can mix only matrix copies within the same central
embedding of the same simple isogeny factor. Such copies, and
opposite embeddings in the dual case, have already been identified.
The proposed link is impossible.

Every self-paired column is mixed in every conjugate, and every
non-CM quadratic-center column is mixed at its two embeddings.
This argument also covers higher-center columns when all their
embeddings are mixed. It does not assert independence of their
central characters.

\subsection{Partially pure Type IV columns}

A partially pure factor has $r_B=1$ and is one of
$(f,a,g)=(2,2,4),(2,3,6),(3,2,6)$.
Choose the representatives of its opposite pairs to be the
embeddings of $K$ lying over one embedding of $E$.
The action of $\Gal(\kk/E)$ on this set is transitive, whether or
not $K/E$ is Galois. Its size $f$ is prime.

An invariant equivalence relation on a prime-sized transitive set
is either discrete or universal: its classes have equal size
dividing $f$. On a column with Hodge multiplicities $(p,a-p)$,
the derived part of the normalized circle differential is
\[
 h'-\frac{2p-a}{a}\operatorname{id}.
 \tag{10}
\]
It is zero exactly when $p=0$ or $a$, namely at a pure column.
It is nonzero at a mixed column. A graph isomorphism identifies
projections of the same source differential and preserves whether
they are zero. Since both kinds occur, the universal relation is
impossible. Therefore the ideals within this factor are independent.

Only the quartic-center fourfold can have an external non-CM
partner: the other two cases already use dimension six.
Its two $A_1$ targets are interchanged by Galois. The remaining
dimension is at most two and inherits $E$. The preceding complete
list shows that a non-CM target in that remainder can only be one
$A_1$: a non-CM elliptic curve repeated twice, or a Type II
surface. This single absolute target is Galois-stable.
A link to one member of the interchanged pair would, by conjugation
and transitivity of linkage, link it to the other, contradicting
their established independence. CM partners contribute no derived
target.

\subsection{The orthogonal exceptions}

For $\SO_6\simeq\SL_4/\{\pm1\}$, the six-dimensional module
$\Lambda^2\kk^4$ differs from the four-dimensional standard
$\SL_4$ module. The associated Type III factor has $g=6$ and
uses the entire dimension budget. There is no second column block
or non-CM factor with which to link it. This disposes of the only
unequal-module simple-target identification in this range outside
$\SO_4$.

For $\SO_4$, the column is $\kk^2\otimes\kk^2$.
Write the two semisimple circle projections as $x_1,x_2$.
If both are nonzero, their standard spectra are $\{a,-a\}$ and
$\{b,-b\}$, $a,b\ne0$, and the tensor spectrum
$\{\pm a\pm b\}$ has at least three elements. Thus at most one
projection is nonzero. The self-pair forces mixedness, so exactly
one is nonzero. The connected center acts trivially on this
self-paired column, so there is no hidden scalar contribution.

These are the only two absolute ideals of the rational derived
group of this simple Type III fourfold: its sole column, repeated
twice, is faithful. If each ideal were separately Galois-stable,
the one with zero differential would give a proper rational
adjoint kernel containing the circle, contradicting minimality.
Galois therefore interchanges the ideals. Their internal
independence is already part of the full $\SO_4$ projection.
An external non-CM remainder again has at most one Galois-stable
$A_1$ target, so the same conjugation argument excludes all
external links.

\begin{proof}[Completion of Theorem~\ref{thm:main}]
All individual projections and all possible graph links have
been treated, so Lemma~\ref{lem:graphs} gives the full product
Lie algebra. The image of the connected algebraic group $R_\kk$
is a connected closed subgroup of the product of the connected
classical matrix groups. Equal Lie algebras imply equal dimension;
a closed subgroup of full dimension in that connected product
is the whole product. Thus this is equality of matrix images,
not merely an isogeny of abstract Lie groups.

On the CM summand the projection of $J$ is toric, so its derived
projection is trivial. Matrix units and opposite pairings determine
the action on all repeated and opposite copies. The forms and maps
are admissible by Lemma~\ref{lem:columns}. This gives exactly
(1) on the original Hodge structure and condition (c) in the
manuscript's sense.
\end{proof}

\section{What the executable certificate proves}

Run, without external packages or a network connection,
\begin{verbatim}
python3 research/actual_group_audit/verify.py
\end{verbatim}
The program compares freshly recomputed data with
\path{exact_certificate.json}. It raises an error on
any mismatch; no timeout, exception, unknown case or empty search
is converted to success. It does not import the source inventories
or the manuscript's claimed representation table as an answer.

\subsection{Coverage, not a sample}

The parameter loop uses the divisibilities in (5), $gr_B\leq6$,
the even-copy field-action restrictions, and the proved CM/rank-two
exclusions. For each Type IV case it enumerates every tuple
$p_\tau\in\{0,\ldots,a\}$ and retains every one not excluded by
Lemma~\ref{lem:pure}. It also records all ranks $r_+,r_-$ for a
split noncentral quadratic action. No particular number field,
period, polarization matrix or tensor degree is sampled.
The result has 68 covering templates: 7 of Type I, 7 of Type II,
2 of Type III, and 52 of Type IV.

All multisets of these templates of total dimension at most six
are enumerated, including the empty non-CM multiset and an arbitrary
CM remainder. Repeated numerical templates are allowed to represent
distinct nonisogenous factors. There are 111 such multisets.
Each has at most three absolute simple targets after splitting
$\mathfrak{so}_4$. Every set partition of these targets is tested,
not merely pairwise links or selected diagrams.

For the partial orbits the code imposes only the necessary cyclic
Galois permutation. In degree two this is the transposition. A
transitive subgroup of $S_3$ contains a three-cycle (otherwise two
distinct transpositions would already generate one, while a single
transposition is not transitive). Thus the cyclic test covers both
$C_3$ and $S_3$, rather than assuming $K/E$ is cyclic. Any external
target is fixed because there is at most one and it is the rational
$A_1$ established above. For $\SO_4$ the required transposition
was proved from circle minimality.

Across the 306 tested partitions, the 111 discrete partitions are
the only survivors. The 195 nontrivial partitions are excluded as
follows (the first applicable obstruction is recorded):
\begin{center}
\begin{tabular}{lr}
Nonisomorphic simple ideals &14\\
Failure of the already full individual $\SO_4$ projection &5\\
Failure of Galois invariance &74\\
Zero/nonzero circle-differential conflict &46\\
All-conjugate intertwiner enlarging the actual commutant &56
\end{tabular}
\end{center}
An arbitrary graph automorphism is not being sampled: table (8)
handles all twists of an ordinary standard module up to duality,
and the zero/nonzero tests and Galois arguments are invariant under
every Lie algebra isomorphism. The preceding proofs establish why
these combinatorial obstructions are necessary for the actual group.
That bridge, not the numerical value zero alone, makes the result
universal.

\subsection{Adversarial controls and the Jacobian boundary}

The diagonal $A_1$ inside the tensor $A_1+A_1$ has image dimension
three and a two-dimensional commutant. It is explicitly computed
as an excluded diagram, not promoted to an abelian-variety
counterexample. The script also compares diagonal and independent
$A_1$ actions on two standard copies: their commutants have
dimensions four and two. The diagonal case is the matrix-copy
phenomenon, for example on the square of a non-CM elliptic curve.
It is removed by using the actual commutant; it does not refute
the lemma after the prescribed identifications.
Likewise the alternating determinant on a pure rank-two column has
circle weight $+2$ or $-2$, whereas on a mixed column it has weight zero.
Both pure orientations are checked.
This exactly detects why derived $\SL_2$ self-duality must not be
silently treated as a Hodge self-pair.

Jacobian was used only through the approved Argus adapter.
The inspected operation was \texttt{matrix.nullspace.compute},
server version \texttt{0.14.0}, protocol \texttt{2025-11-25}.
Its exact payloads and typed replies are retained in the dedicated
JSON files. The full $\SO_4$ invariant-form system has 48 rows
and 16 columns. Although the exposed entry-array schema allowed
64 rows, the runtime rejected this request with
\texttt{matrix.budget\_exceeded}, reporting a 32-row/column limit;
the adapter exited with code 2. This was an interface rejection,
not mathematical evidence.

The equations were replaced by their exact 15-row RREF basis,
with row-space equivalence checked by the standalone program.
The corrected 15-by-16 request exited with code 0 and returned
rank 15 and nullity one, with a symmetric nondegenerate invariant
form. The standalone rational elimination independently recomputes
the full-system kernel and agrees with every entry of the returned
basis. The discovery and inspected contract do not supply a
Hodge-group classification, and no such capability is attributed
to the adapter.

\paragraph{Exactness and limitations.}
The checker verifies integer root reflection closure, rational
Weyl dimensions, Weyl orbits, invariant forms and commutants,
the exceptional exterior-square bracket identities, numerical
parameter coverage, and all finite graph partitions. The analytic
and rational Hodge-theoretic bridge is the proof above, not a
Lean/kernel formalization. No \texttt{.lean} file or compiler
certificate is claimed. Independent review must assess that bridge,
the hypotheses of the four imported inputs, and the correspondence
between the mathematical proof and the enumeration.

\section{Consequences and remaining campaign frontier}

The target source is
\begin{quote}\small\ttfamily
research/openai\_manuscript/\\
Weil-classes-and-Hodge-classes-on-abelian-powers-October-6-2026/\\
build/paper.tex.
\end{quote}
The verification covers the classification and independence chain
at lines 1128--1608, specifically:
\begin{quote}\small\ttfamily
p3ab:albert-columns; p3ab:minuscule;\\
p3ab:small-projections; p3ab:small-columns;\\
p3ab:independence.
\end{quote}
It supplies condition (c) of
\texttt{p3:tensor-criterion} using the actual period and actual
endomorphism algebra.

The proofs of \texttt{p3ab:quadratic-five} at lines 1672--1722 and
\texttt{p3ab:split-sixfolds} at lines 1724--1834 both invoke this
independence step. This removes one obstruction to their
all-powers deductions. It does \emph{not} verify their simultaneous
balance equation, the CM supply of the required determinant
products, the invariant-theoretic tensor criterion as a whole,
the external CM Hodge theorem, or the cycle construction.
In particular we do not infer the rational Hodge conjecture from
Theorem~\ref{thm:main}.

The geometric theorem \texttt{p3sw:weil-plane} is on a different
branch. Its final proof, lines 8004--8024, invokes
\texttt{p3sw:endpoint-all-lattices}; the inventoried geometric
branch does not invoke \texttt{p3ab:independence}. Its dependence
on the earlier Hermitian convention and the elementary
hyperbolic-basis statement must not be confused with dependence
on the actual-group argument. No direct source-declared dependence
of the staged Kuga--Satake or main Milne rationality theorems on
this independence lemma was established by the complete source
reading.

All 21 staged \texttt{.tex} files were read, and 191 source claim
records, including 130 formal theorem/lemma/proposition/corollary
statements, were added to the existing shared ledger and proof
graph. Manuscript-specific hypotheses are retained; those source
assertions are not marked verified merely because they were
catalogued. The seven \texttt{AG-} audit claims separate the new
proof obligations from that source inventory. The generic-group
substitution route is retired with its reason in the ledger.
The source files themselves are unchanged.

This bounded mission submits a positive verification theorem, not
an exact counterexample. There is no surviving hypothesis-valid
obstruction in the proved reduction. Publication-level novelty and
independent certification are separate requirements, and a
known-result reconstruction or finite exclusion alone would not
close the campaign.

\begin{thebibliography}{9}
\bibitem{CO}
C.-L. Chai and F. Oort,
\emph{Moduli of abelian varieties and $p$-divisible groups},
in \emph{Arithmetic Geometry}, Clay Mathematics Proceedings 8 (2009),
441--536. ISBN 978-0-8218-4476-2.
\url{https://www.claymath.org/wp-content/uploads/2022/03/cmip08c.pdf}.
Locators: Sections 10.12--10.14.

\bibitem{MZ}
B. J. J. Moonen and Yu. G. Zarhin,
\emph{Hodge classes on abelian varieties of low dimension},
Mathematische Annalen 315 (1999), 711--733.
The version used for all numbered locators is
\href{https://arxiv.org/abs/math/9901113v2}{arXiv:math/9901113v2},
13 April 1999. Locators: Theorems (0.1)--(0.2),
Sections (1.1)--(1.3), Section (3.1), Lemmas (3.3)--(3.4).

\bibitem{Et}
P. Etingof, \emph{Lie groups and Lie algebras},
\href{https://arxiv.org/abs/2201.09397v4}{arXiv:2201.09397v4} (2024).
Locators: Example 23.5, Theorem 23.7, Remark 23.18,
Proposition 26.8, Section 27.2, Definition 30.1, Lemma 30.2,
Corollary 30.5, Proposition 30.6, Section 30.4.
\end{thebibliography}
\end{document}
